\documentclass{ar2rc}

\usepackage{commands}
\addbibresource{ref.bib}
\setmainpaper{demo/two-column/manuscript}

\papertitle{Selecting Latency Budgets for Edge Inference}
\manusid{ACM-DEMO-TWO-COLUMN}
\authors{San Zhang}
\journalname{ACM Review Demonstration}

\begin{document}
\responseletterheader

Thank you for reviewing the manuscript. The badges link to revisions in the
appended two-column manuscript, and the margin badges return to the replies.

\responsesection{Responses to Reviewer 1}
\begin{reviewerblock}
\begin{reviewercomment}
Specify when the latency measurement starts and give the time budget.
\end{reviewercomment}
\begin{reviewerresponse}
We now define the measurement interval and state the budget used for all
devices. See \elinks[R1-C1]{E3,E1}.
\end{reviewerresponse}
\end{reviewerblock}

\responsesection{Responses to Reviewer 2}
\begin{reviewerblock}
\begin{reviewercomment}
Clarify whether the budget changes between workloads.
\end{reviewercomment}
\begin{reviewerresponse}
We now state explicitly that both workloads use the same budget. See
\elink[R2-C1]{E2}.
\end{reviewerresponse}
\end{reviewerblock}

\begin{reviewerblock}
\begin{reviewercomment}
Explain how the baseline is calibrated.
\end{reviewercomment}
\begin{reviewerresponse}
We now describe the separate calibration trace. See \elink[R2-C2]{E4}.
\end{reviewerresponse}
\end{reviewerblock}

\appendmainpaper
\end{document}
